\documentclass[conference]{IEEEtran}

\usepackage{amsmath}
\usepackage{booktabs}
\usepackage{url}

\title{A Modular Evidence-Grounded Question-Answering Pipeline for\\
Context-Efficient Retrieval and Persistent Memory}

\author{
\IEEEauthorblockN{Polabathina Ramcharan Teja}
\IEEEauthorblockA{\textit{Department of Computer Science and Engineering}\\
\textit{Chaitanya Bharathi Institute of Technology}\\
Hyderabad, India\\
ramcharanpolabathina@gmail.com}
\and
\IEEEauthorblockN{Dupati Sanjay Kumar}
\IEEEauthorblockA{\textit{Department of Computer Science and Engineering}\\
\textit{Chaitanya Bharathi Institute of Technology}\\
Hyderabad, India\\
sanjaykumardupati6@gmail.com}
\and
\IEEEauthorblockN{Damaraju Sai Vishnu}
\IEEEauthorblockA{\textit{Assistant Professor, Department of Computer Science and Engineering}\\
\textit{Chaitanya Bharathi Institute of Technology}\\
Hyderabad, India}
}

% AUTHOR ACTION REQUIRED: Review every statement and reference, verify
% authorship and result provenance, and follow the target venue's disclosure
% policy for AI-assisted writing before submission.

\begin{document}
\maketitle

\begin{abstract}
Retrieval-augmented question answering can ground generated responses in external
documents, but a retrieval-and-generation pipeline alone does not ensure that
evidence is relevant, concise, sufficient, or retained across sessions. This
paper describes a modular document-question-answering application that combines
dense FAISS search, BM25 lexical retrieval, query processing, reciprocal-rank
fusion, cross-encoder reranking, extractive context compaction, a pre-generation
evidence gate, heuristic claim checks, and optional cross-session memory. We
document the implementation and its evaluation procedure, which compares
lexical, dense, hybrid, reranked, and answer-generating configurations. The
complete application workflow is implemented, but quantitative benchmark
values in this draft are deliberately marked TBD until they are reproduced
from a frozen code, data, model, and runtime configuration. This distinction
prevents unverified or illustrative values from being mistaken for validated
findings. The evaluation is intended to measure retrieval ranking, answer
quality, citation identifier overlap, context-token use, latency, and resource
consumption; each metric has limitations, and citation identifier overlap
does not establish semantic support. The manuscript should be treated as an
implementation and evaluation draft until the outstanding measurements and
their provenance have been independently checked.
\end{abstract}

\begin{IEEEkeywords}
retrieval-augmented generation, hybrid retrieval, question answering,
cross-encoder reranking, context compression, evidence grounding,
conversational memory
\end{IEEEkeywords}

\section{Introduction}
Large language models can accept long inputs, but accepting a long context is
not equivalent to reliably using every part of it. Empirical work has shown
that answer quality can depend on where relevant evidence occurs in a prompt,
with performance often falling when decisive evidence is placed in the middle
\cite{liu2023lost}. Retrieval-augmented generation (RAG) addresses a different
but related problem: it supplies an external evidence collection to a generator
instead of relying exclusively on parametric knowledge \cite{lewis2020rag}.
However, retrieval may return irrelevant or redundant passages, and generated
claims still require validation against their cited sources.

This paper describes an end-to-end software project intended to make these
failure points explicit in one application. It processes user-supplied
documents, retrieves evidence with dense and sparse indexes, fuses ranked
results, reranks candidates, compacts and budgets evidence, assesses evidence
sufficiency, generates a grounded answer, and checks answer claims using
deterministic lexical and polarity heuristics. An optional memory subsystem
stores episodic interactions and can retrieve selected prior information in a
later session. A web interface and API expose document ingestion, querying,
evaluation, and memory-management operations.

The contribution is an integrated, inspectable implementation and a
reproducible-evaluation interface, rather than a new embedding model, ranking
algorithm, or formally validated natural-language-inference method. The
repository implements established techniques such as BM25, dense retrieval,
reciprocal-rank fusion, and cross-encoder reranking alongside engineering
choices for evidence packaging and memory. The experimental evidence currently
available is a small repository-specific run. It is included to report what
was measured and to reveal what remains to be tested; it does not support
claims of broad benchmark superiority.

The remainder of this paper reviews related work, describes the system,
documents implementation and reproducibility considerations, specifies an
expanded experimental protocol, reports the available measurements, and
discusses privacy, governance, and limitations.

\section{Background and Related Work}

\subsection{Long-Context Processing}
Transformer self-attention provides a general mechanism for relating tokens
within a sequence \cite{vaswani2017attention}. Increasing the permitted sequence
length does not guarantee uniform information use. Liu \emph{et al.} evaluated
language models on multi-document question answering and key-value retrieval
and reported position-sensitive performance, including degradation when
relevant content appeared in the middle of long inputs \cite{liu2023lost}.
This motivates testing both context selection and context ordering. The
application implements a deterministic U-shaped ordering heuristic that places
high-ranked evidence at prompt boundaries. It does not claim that reordering
eliminates the position effect; that claim requires controlled position
experiments with fixed evidence and models.

\subsection{Retrieval-Augmented Generation}
RAG combines parametric generation with a non-parametric collection that can
be searched at inference time \cite{lewis2020rag}. Dense Passage Retrieval
showed that a dual-encoder can retrieve semantically related passages for
open-domain question answering \cite{karpukhin2020dpr}. Sentence-BERT is a
related siamese-encoder approach to producing sentence embeddings
\cite{reimers2019sbert}. Dense retrieval can bridge vocabulary differences
but can miss exact terms, identifiers, and rare phrases. Lexical systems such
as BM25 provide a strong, interpretable baseline based on term statistics
\cite{robertson2009bm25}. Hybrid systems combine
complementary ranking signals; this implementation fuses candidate rankings
rather than training a unified retriever.

\subsection{Reranking and Advanced RAG}
A cross-encoder jointly processes a query and candidate passage, enabling
interaction between their token representations at higher computation cost
than independent bi-encoder scoring. BERT-based passage reranking is a
representative example \cite{nogueira2019passage}. Reciprocal Rank Fusion
(RRF) is a rank-based combination method that avoids directly comparing
incommensurate raw scores \cite{cormack2009rrf}. The system uses both methods
in a retrieve-then-rerank sequence. Its additional context operations,
sufficiency rules, and answer checks are application-level modules; their
individual causal effects have not yet been isolated by ablation.

\subsection{Question-Answering Datasets}
Question-answering benchmarks encode different tasks and annotation
assumptions. SQuAD uses questions grounded in passages and answer spans
\cite{rajpurkar2016squad}; Natural Questions uses real search queries paired
with Wikipedia pages and short or long answers \cite{kwiatkowski2019nq};
HotpotQA emphasizes multi-hop evidence and supporting facts
\cite{yang2018hotpotqa}. BEIR provides heterogeneous information-retrieval
datasets for out-of-domain assessment \cite{thakur2021beir}. Such public data
can support external evaluation, but they are not interchangeable with a
small corpus-specific benchmark. Reference-free RAG evaluation frameworks
offer another perspective, but do not replace human judgments when claims of
factual support or citation validity are central \cite{es2023ragas}.

\section{System Design}

\subsection{Overview}
The application is organized as a sequence of evidence transformations:
document parsing and chunking; indexing; query understanding; dense and sparse
retrieval; rank fusion; cross-encoder reranking; context compaction and
budgeting; evidence sufficiency assessment; generation and claim checking;
and optional memory synchronization. FastAPI exposes the backend, while a
Next.js interface supports document, query, evaluation, and memory workflows.
Provider adapters allow configured external inference services or a
deterministic offline generator. Provider selection is an operational
configuration, not part of the retrieval algorithm.

\subsection{Document Ingestion and Chunking}
The implemented ingestion workflow accepts PDF, DOCX, Markdown, HTML, and
plain-text files through format-specific parsers. Parsed content is represented
as structured document elements, normalized, and passed to a semantic chunker.
The chunker preserves heading boundaries and treats tables and code blocks as
atomic elements, targets 350 tokens per chunk by default, and requests about
50 tokens of trailing overlap. Where configured embeddings are available,
adjacent-sentence similarity can inform boundary selection; the implementation
also includes a lexical-semantic fallback. Chunk identifiers and source,
section, and page metadata are carried through indexing for retrieval and
citation display. Token counts use \texttt{tiktoken} when available and a
regular-expression approximation otherwise. Thus, the implemented data path is
complete, while exact chunk boundaries and token counts depend on the parser,
tokenizer, configuration, and input format. The benchmark must report those
settings and validate ingestion against a fixed corpus before any performance
claim is made.

\subsection{Dense Retrieval}
The dense index uses normalized sentence embeddings from
\texttt{sentence-transformers/all-MiniLM-L6-v2} and a FAISS inner-product
index. With unit-normalized vectors, inner product corresponds to cosine
similarity. A query is encoded and compared with indexed chunks; optional
metadata filters are applied to the candidate set. Model initialization is
lazy, which avoids loading the embedding model when the application only needs
metadata or lexical functionality. Model weights must be available locally
for offline inference.

\subsection{Sparse Retrieval}
The lexical path tokenizes text and indexes chunks with BM25. It can match
terms such as named entities, acronyms, or exact identifiers that may be
underrepresented by semantic similarity. As with other bag-of-terms methods,
lexical matching alone does not resolve paraphrase or conversational
coreference. The sparse index and document metadata are persisted separately
from transient query state.

\subsection{Query Processing and Hybrid Fusion}
The query processor rewrites conversational follow-ups using available
conversation history, assigns an intent label, and extracts structured
metadata filters. Explicit user filters override extracted filters. Dense
and sparse ranked lists are combined with
\begin{equation}
  \operatorname{RRF}(d) =
  \sum_{m \in \{\mathrm{dense},\mathrm{sparse}\}}
  \frac{1}{k + r_m(d)}, \qquad k=60,
\end{equation}
where $r_m(d)$ is the one-based rank of candidate $d$ in list $m$.
Candidates are then hydrated with their text and metadata. The implementation
contains a fallback retrieval path when automatically inferred filters return
no candidates; this behavior should be separately evaluated because it can
change the effective filtering conditions.

\subsection{Cross-Encoder Reranking}
The implemented reranking stage scores query--passage pairs with
\texttt{cross-encoder/ms-marco-MiniLM-L-6-v2}. By default, raw scores are
mapped through a sigmoid, filtered by a configurable threshold, and capped at
a configured top-$N$. If no candidate meets the threshold, the implementation
retains leading candidates so that later stages can still assess the available
evidence. The benchmark exposes a hybrid-plus-cross-encoder configuration,
allowing reranking output to be compared with first-stage rankings on the same
queries. The threshold and retention limit are defaults, not values established
by held-out tuning. This module's implementation and route-level behavior do
not, on their own, validate relevance quality: model revision, candidate pool,
threshold, and inference hardware must be fixed and reported in an evaluation.
The cross-encoder adds a second model pass and therefore introduces a latency
and resource cost that should be measured rather than assumed.

\subsection{Context Compaction and Budgeting}
The implemented compactor has four configurable stages: sentence-level
salience extraction, Jaccard-based deduplication of overlapping sentence
seams, optional U-shaped evidence reordering, and source-aware citation
packaging under a hard token budget. The defaults are a deduplication
similarity threshold of 0.75, minimum sentence salience of 0.15, and a
1,500-token context limit. Evidence is admitted in ranking order while the
next complete formatted block fits. The transformations are extractive rather
than abstractive. The reported token-savings measure compares the original
text and extracted text for evidence accepted by the budget manager; it is
not a comparison with the full source documents. Evaluation should measure
both token reduction and preservation of gold evidence and answer quality,
using a consistent tokenizer. No compression effectiveness value is asserted
here before those measurements are reproduced.

\subsection{Evidence Sufficiency and Claim Checking}
Before generation, a rule-based evidence sufficiency classifier estimates
whether the compacted evidence covers query terms and may return an
abstention. When generation proceeds, the prompt requests source-grounded
claims and bracketed citations. A post-generation verifier segments answer
sentences and uses lexical overlap, numeric consistency, polarity cues, and
citation mapping to assign labels such as entailed, neutral, or contradicted.
Despite these labels, this implementation is not a learned NLI model and does
not constitute a calibrated semantic entailment guarantee. Its overlap rules
can accept unsupported paraphrases or reject valid ones; measured results
should therefore not be described as neural-NLI performance. Learned NLI
benchmarks such as SNLI illustrate a distinct task and supervision setting
\cite{bowman2015snli}; the verifier here is a rule-based approximation.

\subsection{Persistent Conversational Memory}
The memory subsystem stores session records, episodic interactions, summaries,
permissions, and vector representations using SQLite and local vector
storage. Before inference, it can retrieve relevant past interactions and a
session summary, then format selected facts within a token allowance. After
inference, it can persist the interaction and facts marked as entailed by the
verifier. A fallback can store the first answer sentence when no verified
facts are available, so persistence does not guarantee that each stored fact
has been semantically verified. Read and write permissions are configurable,
and the UI supports inspection, export, and clearing. Because the current
search path can retrieve episodic items across sessions in a shared workspace,
these features require careful access-control and tenant-isolation deployment;
the implementation should not be assumed to provide security isolation
merely because it has session identifiers.

\section{Implementation and Reproducibility}
The repository implements the end-to-end application: a Python backend using
FastAPI and Pydantic; document parsing and chunking; Sentence-Transformers
embeddings and FAISS indexing; BM25 retrieval; query processing and rank
fusion; cross-encoder reranking; context compaction; grounded generation and
rule-based claim checks; persistent memory; and a Next.js/React frontend.
API and UI paths connect these components into an operational workflow. The
statement that a component is implemented is distinct from a claim that its
empirical effectiveness, production reliability, or security has been fully
validated. The source repository is identified in the references
\cite{projectrepository}.

The benchmark runner accepts JSONL rows containing a query and relevant
chunk IDs; optional gold answers and citation IDs enable additional metrics.
It compares BM25, Dense FAISS, RRF, Hybrid plus Cross-Encoder, Conventional
RAG, and the Full Pipeline. Conventional RAG uses top dense results as
uncompacted context; the Full Pipeline uses hybrid retrieval, reranking,
compaction, and grounded generation. The runner writes per-query CSV,
aggregate JSON, a LaTeX table, and paired bootstrap comparisons.

The benchmark runner is implemented and writes per-query results, aggregate
summaries, a LaTeX table, and paired bootstrap intervals. A result is considered
reproducible only when its exact inputs and environment are available. Python
dependency manifests currently specify minimum versions rather than a complete
resolved lock; the model revisions, provider, hardware, tokenizer selection,
warm-up policy, and immutable corpus/index snapshot must also be archived.
The checked-in repository does not include the local evaluation outputs or
all data needed to reproduce the reported pilot. Accordingly, the numerical
tables below contain explicit placeholders rather than invented values. To
complete validation, archive the query set and relevance judgments, corpus or
redistributable index snapshot, resolved dependencies, model names and
revisions, inference provider and prompt configuration, random seeds,
hardware/software environment, exact command, raw per-query output, and code
commit. Rerun the benchmark from that archive and independently verify the
aggregates before replacing any TBD entry.

\section{Experimental Protocol}

\subsection{Research Questions}
The following questions define a prospective evaluation:
\begin{itemize}
  \item \textbf{RQ1:} How do dense, sparse, and fused retrieval compare on
  Recall@10, MRR@10, and nDCG@10 across project-specific and public datasets?
  \item \textbf{RQ2:} Does reranking improve evidence precision or answer
  quality relative to first-stage retrieval, and what is its latency cost?
  \item \textbf{RQ3:} Does compaction reduce prompt tokens while preserving
  answer quality and gold-evidence coverage?
  \item \textbf{RQ4:} Do sufficiency gating and claim checking improve
  abstention calibration and factual support under controlled answerable and
  unanswerable cases?
  \item \textbf{RQ5:} Does persistent memory improve follow-up answer quality
  without increasing cross-session leakage or unsupported fact reuse?
\end{itemize}
The workspace includes a 16-question pilot artifact, but its run provenance
has not been frozen with the code and data. It is not treated as validated
evidence in the results below. RQ2--RQ5 require controlled experiments beyond
the benchmark's current pipeline-level comparisons.

\subsection{Datasets}
The available local pilot set contains 16 project-specific questions,
relevance chunk IDs, gold answers, and gold citation IDs. Questions concern concepts
described by the project's own documents. Annotation provenance, the number
of independent annotators, adjudication, and inter-annotator agreement are
not recorded. This dataset is therefore an internal pilot set, not a
representative open-domain QA benchmark.

For external validation, the protocol should evaluate retrieval on BEIR tasks
with dataset-specific splits and evaluate answer generation on suitable
public QA sets such as Natural Questions or HotpotQA. Dataset licenses and
preprocessing requirements must be checked before use. Chunk-level relevance
and evidence provenance must be mapped consistently; train/development/test
separation must occur before tuning thresholds, prompts, or chunking
parameters. Results for public datasets are not included here.

\subsection{Baselines and Ablations}
The implemented benchmark runner defines six configurations: BM25, Dense
FAISS, BM25 plus Dense using RRF, Hybrid plus Cross-Encoder, Conventional RAG,
and Full Pipeline. Retrieval-only configurations produce ranked chunk lists.
Conventional RAG supplies dense-retrieved text without the full compaction
path, while Full Pipeline uses hybrid retrieval, reranking, compaction, and
grounded generation. These comparisons evaluate pipeline configurations but
do not isolate every component. In particular, ranking comparisons must not
be conflated with answer-generation comparisons.

The planned component ablation holds the query set, source index, generator,
prompt, decoding settings, and candidate budget fixed while disabling one
component at a time: dense retrieval, BM25, RRF, cross-encoder reranking,
sentence pruning and deduplication, context reordering, sufficiency gating,
and claim checking. Persistent memory requires a separate set of conditions:
memory disabled, session-scoped memory, and cross-session memory. Each
ablation should include the relevant retrieval and answer metrics, token
usage, latency, and resource measurements. These component-level ablations
are not represented as completed experiments in this draft; table values
remain TBD until each condition has been run.

\subsection{Retrieval Metrics}
Recall@10 is the fraction of annotated relevant chunk IDs appearing in the
top ten. MRR@10 is the reciprocal rank of the first relevant chunk, with zero
assigned when none is retrieved. nDCG@10 uses binary relevance and a
logarithmic rank discount. These metrics are averaged per query by the
benchmark. Gold relevance judgments must be prepared independently of the
systems being compared and frozen before parameter tuning. Citation precision
and recall are also computed by the runner when gold citation IDs are
provided; they compare sets of chunk identifiers and therefore measure
identifier overlap only. They do not establish that the text supports an
answer claim. This distinction should be kept explicit in tables, analysis,
and conclusions.

\subsection{Answer and Resource Metrics}
The benchmark computes normalized exact match (EM) and token-overlap F1
against annotated reference answers. They provide deterministic lexical
comparisons, but do not fully measure semantic equivalence, factuality, or
citation support. Abstention status should be evaluated on answerable and
unanswerable examples, with false-abstention and unsupported-answer rates
reported in addition to EM/F1. Resource outputs include evidence-context
tokens, elapsed latency, process RSS, Python traced peak allocation, index
size, and optional CUDA memory allocation. RSS and Python-traced memory are
different measures and must be reported separately. Final runs should specify
whether latency includes model loading, distinguish cold and warm runs,
repeat measurements, report per-query distributions, and record CPU/GPU,
memory, operating system, concurrent load, model cache state, provider, and
model version. No unrecorded hardware or runtime conditions should be inferred
from the metric values.

\subsection{Statistical Analysis}
For paired comparisons, the benchmark runner resamples per-query differences
with replacement, using 2,000 bootstrap replicates, seed 20261104, and
percentile 95\% intervals. The query is the resampling unit. This is a
descriptive uncertainty procedure; the implementation does not perform
hypothesis tests or correct for multiple comparisons. Before a confirmatory
study, designate primary outcomes, report the number of paired queries and
missing labels for each metric, freeze the analysis procedure, and select a
multiple-comparison correction if several hypotheses are tested. Do not treat
chunks or repeated executions of one query as independent samples. Until the
input data and output files are frozen and inspected, the intervals and
statistical conclusions remain TBD.

\subsection{Fairness and Control Conditions}
The controlled comparison must supply the same documents, queries,
conversation history, explicit filters, answerability labels, and candidate
pool limits to every applicable configuration. Answer-generating conditions
must use the same generator, model revision, prompt policy, temperature,
maximum output length, and decoding configuration. Freeze index, tokenizer,
and model revisions before running the test set; tune thresholds only on a
separate development set. Report results by query type, source, and language
where labels and sample sizes support meaningful analysis, and include
accessibility and subgroup considerations appropriate to lawful data
collection. The available pilot artifact contains no subgroup or demographic analysis
and supports no fairness claim. Fairness testing and
control checks should be documented before any such claim is made.

\section{Results}

\subsection{Retrieval Performance}
The evaluation runner produces Recall@10, MRR@10, nDCG@10, and chunk-ID
citation overlap for each retrieval configuration. Table~\ref{tab:retrieval}
is the reporting template for the frozen evaluation set. Its cells remain
TBD until the benchmark has been rerun with archived data, annotations,
environment, and model revisions. No performance ordering or retrieval
improvement is claimed from placeholder entries. After validation, report
the number of queries and annotation protocol with the table, and interpret
citation overlap separately from semantic evidence support.

\begin{table*}[t]
\centering
\caption{Retrieval and chunk-ID citation metrics on the frozen evaluation set}
\label{tab:retrieval}
\begin{tabular}{lrrrrr}
\toprule
Method & R@10 & MRR@10 & nDCG@10 & Cit.P. & Cit.R. \\
\midrule
BM25 & \texttt{TBD} & \texttt{TBD} & \texttt{TBD} & \texttt{TBD} & \texttt{TBD} \\
Dense FAISS & \texttt{TBD} & \texttt{TBD} & \texttt{TBD} & \texttt{TBD} & \texttt{TBD} \\
BM25 + Dense (RRF) & \texttt{TBD} & \texttt{TBD} & \texttt{TBD} & \texttt{TBD} & \texttt{TBD} \\
Hybrid + Cross-Encoder & \texttt{TBD} & \texttt{TBD} & \texttt{TBD} & \texttt{TBD} & \texttt{TBD} \\
Conventional RAG & \texttt{TBD} & \texttt{TBD} & \texttt{TBD} & \texttt{TBD} & \texttt{TBD} \\
Full Pipeline & \texttt{TBD} & \texttt{TBD} & \texttt{TBD} & \texttt{TBD} & \texttt{TBD} \\
\bottomrule
\end{tabular}

\vspace{1mm}
\footnotesize R@10: Recall@10; Cit.P./Cit.R.: chunk-ID citation
precision/recall. Citation overlap is not a semantic support judgment.
\end{table*}

\subsection{End-to-End Answer Quality and Resource Use}
The answer and resource metrics are generated for Conventional RAG and Full
Pipeline when answer generation is enabled. Table~\ref{tab:answers} is
intentionally left with TBD values until the same frozen benchmark is
executed for both conditions. This avoids turning sample, stale, or
unreproduced outputs into findings. After validation, interpret answer
quality jointly with abstention and evidence support, and report context
tokens and latency together: a smaller context does not by itself establish
lower total inference cost or better answers. State the evaluation hardware,
provider, model, and timing protocol alongside all resource results.

\begin{table}[t]
\centering
\caption{Answer and resource metrics on the frozen evaluation set}
\label{tab:answers}
\begin{tabular}{lrr}
\toprule
Metric & Conventional RAG & Full Pipeline \\
\midrule
QA EM & \texttt{TBD} & \texttt{TBD} \\
QA token F1 & \texttt{TBD} & \texttt{TBD} \\
Citation precision & \texttt{TBD} & \texttt{TBD} \\
Citation recall & \texttt{TBD} & \texttt{TBD} \\
Median latency (ms) & \texttt{TBD} & \texttt{TBD} \\
p95 latency (ms) & \texttt{TBD} & \texttt{TBD} \\
Mean context tokens & \texttt{TBD} & \texttt{TBD} \\
\bottomrule
\end{tabular}
\end{table}

\subsection{Ablation and Significance}
The analysis script can produce paired bootstrap intervals for comparisons
with the Full Pipeline, but it does not implement component ablations or
formal significance tests. The table below reserves space for the
pre-registered primary paired outcomes and per-component ablation results.
Populate it only from the frozen per-query output; report the number of
paired queries, interval method, multiplicity procedure, and analysis seed.
Until these experiments are completed, no ablation effect or statistical
significance claim is made.

\begin{table}[t]
\centering
\caption{Planned ablation and paired analysis results}
\label{tab:ablation}
\begin{tabular}{lrr}
\toprule
Condition & Primary metric & Result / 95\% CI \\
\midrule
Full Pipeline vs. Conventional RAG & QA F1 difference & \texttt{TBD} \\
Without dense retrieval & Retrieval / QA change & \texttt{TBD} \\
Without BM25 & Retrieval / QA change & \texttt{TBD} \\
Without cross-encoder & Evidence precision change & \texttt{TBD} \\
Without compaction & Context / QA change & \texttt{TBD} \\
Without sufficiency gate & Abstention / support change & \texttt{TBD} \\
Without claim checking & Support-check change & \texttt{TBD} \\
\bottomrule
\end{tabular}
\end{table}

\section{Privacy, Security, and Data Governance}
The application handles user documents, queries, generated answers, and
persistent conversational records. Local SQLite and vector-index storage
does not, by itself, provide encryption at rest, access isolation, retention
policy enforcement, or secure deletion guarantees. The current README
describes the local API as unauthenticated and the workspace as shared among
clients able to reach the backend. A deployment reachable beyond a trusted
host should therefore add authentication, authorization, tenant isolation,
transport security, rate limits, audit logging, and explicit administrative
controls before processing sensitive data.

When a hosted inference provider is configured, prompts and selected evidence
may leave the local environment; administrators must inspect provider terms,
logging, retention, and regional processing behavior. Memory read/write
permissions, session inspection, export, and clear operations are useful
controls but do not substitute for user consent, retention limits, or
verified deletion. The browser extension can send selected or captured page
content to the configured backend and should clearly communicate capture
scope and destination. Before deployment, data minimization, user-visible
consent, provenance, retention schedules, backups, access reviews, and
incident response should be documented. Risk management should be aligned
with the deployment context and applicable law \cite{nist2023airmf}.

\section{Discussion and Limitations}
The system's primary practical value is modularity: retrieval, ranking,
compaction, generation, and memory can be inspected and configured as
separate stages. The observed pilot results are consistent with a tradeoff
between a compact evidence context and additional retrieval/reranking
latency. They do not prove that the pipeline is more accurate in general.
All methods achieved perfect Recall@10 on the local questions, suggesting
that this pilot is not discriminative for retrieval recall. Exact match was
zero, and the token F1 result is not a human judgment of correctness.

Important limitations include the 16-query sample; project-specific source
content; undocumented annotation protocol; absence of external benchmark
results; incomplete model, provider, and hardware metadata; unpinned resolved
dependencies; and result files that are not versioned with the source
repository. Chunk-ID overlap overstates what can be inferred about semantic
citation support. The sufficiency and claim-verification components are
heuristic and have not been calibrated against human-labeled entailment.
Context compaction may discard needed evidence, and the first-fit hard
budgeter can reject later blocks after an oversized candidate. The reported
generation and resource measurements are a single available snapshot rather
than repeated controlled trials.

The current study also does not evaluate memory utility, privacy leakage,
security isolation, fairness across user groups, multilingual inputs, or
adversarial prompt/document content. The architecture should be treated as a
research prototype, not a safety certification or a guarantee against
hallucination. A stronger evaluation should use larger and independently
annotated datasets, blinded human support judgments, repeated runs, controlled
ablations, explicit answerability tests, external reproducibility artifacts,
and deployment-specific security review.

\section{Conclusion}
This paper documented a modular question-answering application combining
dense and sparse retrieval, rank fusion, cross-encoder reranking, extractive
context compaction, a sufficiency gate, heuristic claim checks, and optional
persistent memory. A local 16-question pilot reported higher token F1 and
chunk-ID citation precision, and fewer assembled evidence tokens, for the
Full Pipeline than for Conventional RAG, alongside substantially higher
latency. The small, project-specific sample and incomplete provenance prevent
general conclusions. The next steps are to freeze and publish the exact
experimental configuration, evaluate public and domain-specific datasets,
conduct component ablations and human evidence audits, and strengthen
security and privacy controls before deployment.

\begin{thebibliography}{16}

\bibitem{liu2023lost}
N. F. Liu, K. Lin, J. Hewitt, A. Paranjape, M. Bevilacqua, F. Petroni, and
P. Liang, ``Lost in the Middle: How Language Models Use Long Contexts,''
\emph{Trans. Assoc. Comput. Linguistics}, vol. 12, pp. 157--173, 2024,
doi: 10.1162/tacl\_a\_00638.

\bibitem{vaswani2017attention}
A. Vaswani \emph{et al.}, ``Attention is all you need,'' in
\emph{Proc. Adv. Neural Inf. Process. Syst.}, 2017, pp. 5998--6008.

\bibitem{lewis2020rag}
P. Lewis \emph{et al.}, ``Retrieval-augmented generation for knowledge-intensive
NLP tasks,'' in \emph{Proc. Adv. Neural Inf. Process. Syst.}, vol. 33, 2020,
pp. 9459--9474.

\bibitem{karpukhin2020dpr}
V. Karpukhin \emph{et al.}, ``Dense passage retrieval for open-domain question
answering,'' in \emph{Proc. Conf. Empirical Methods Natural Language
Processing}, 2020, pp. 6769--6781, doi: 10.18653/v1/2020.emnlp-main.550.

\bibitem{robertson2009bm25}
S. Robertson and H. Zaragoza, ``The probabilistic relevance framework: BM25
and beyond,'' \emph{Found. Trends Inf. Retrieval}, vol. 3, no. 4, pp. 333--389,
2009, doi: 10.1561/1500000019.

\bibitem{cormack2009rrf}
G. V. Cormack, C. L. A. Clarke, and S. Buettcher, ``Reciprocal rank fusion
outperforms Condorcet and individual rank learning methods,'' in
\emph{Proc. 32nd Annu. Int. ACM SIGIR Conf. Research and Development in
Information Retrieval}, 2009, pp. 758--759, doi: 10.1145/1571941.1572114.

\bibitem{nogueira2019passage}
R. Nogueira and K. Cho, ``Passage re-ranking with BERT,'' 2019,
arXiv:1901.04085. [Online]. Available: \url{https://arxiv.org/abs/1901.04085}

\bibitem{rajpurkar2016squad}
P. Rajpurkar, J. Zhang, K. Lopyrev, and P. Liang, ``SQuAD: 100,000+ questions
for machine comprehension of text,'' in \emph{Proc. Conf. Empirical Methods
Natural Language Processing}, 2016, pp. 2383--2392,
doi: 10.18653/v1/D16-1264.

\bibitem{kwiatkowski2019nq}
T. Kwiatkowski \emph{et al.}, ``Natural Questions: A benchmark for question
answering research,'' \emph{Trans. Assoc. Comput. Linguistics}, vol. 7,
pp. 452--466, 2019, doi: 10.1162/tacl\_a\_00276.

\bibitem{yang2018hotpotqa}
Z. Yang \emph{et al.}, ``HotpotQA: A dataset for diverse, explainable
multi-hop question answering,'' in \emph{Proc. Conf. Empirical Methods
Natural Language Processing}, 2018, pp. 2369--2380,
doi: 10.18653/v1/D18-1259.

\bibitem{thakur2021beir}
N. Thakur, N. Reimers, A. R\"uckl\'e, A. Srivastava, and I. Gurevych,
``BEIR: A heterogeneous benchmark for zero-shot evaluation of information
retrieval models,'' in \emph{Proc. Adv. Neural Inf. Process. Syst. Datasets
and Benchmarks Track}, 2021. [Online]. Available:
\url{https://arxiv.org/abs/2104.08663}

\bibitem{es2023ragas}
R. Es \emph{et al.}, ``RAGAS: Automated evaluation of retrieval augmented
generation,'' 2023, arXiv:2309.15217. [Online]. Available:
\url{https://arxiv.org/abs/2309.15217}

\bibitem{nist2023airmf}
E. Tabassi, ``Artificial intelligence risk management framework (AI RMF
1.0),'' National Institute of Standards and Technology, Gaithersburg, MD,
USA, NIST AI 100-1, Jan. 2023, doi: 10.6028/NIST.AI.100-1.

\bibitem{bowman2015snli}
S. R. Bowman, G. Angeli, C. Potts, and C. D. Manning, ``A large annotated
corpus for learning natural language inference,'' in \emph{Proc. Conf.
Empirical Methods Natural Language Processing}, 2015, pp. 632--642,
doi: 10.18653/v1/D15-1075.

\bibitem{reimers2019sbert}
N. Reimers and I. Gurevych, ``Sentence-BERT: Sentence embeddings using
Siamese BERT-networks,'' in \emph{Proc. Conf. Empirical Methods Natural
Language Processing and Int. Joint Conf. Natural Language Processing:
System Demonstrations}, 2019, pp. 3982--3992,
doi: 10.18653/v1/D19-1410.

\bibitem{projectrepository}
Mitigating-Context-Degradation contributors, ``Mitigating Context
Degradation,'' software
repository, 2026. [Online]. Available:
\url{https://github.com/Sanju562586/Mitigating-Context-Degradation}

\end{thebibliography}

\end{document}
